\subsection{指数関数と対数関数}
$a>0$ として,指数関数 $a^x$ を定義から微分してみよう.指数法則 $a^{x+h}=a^xa^h$ から
\[
 \frac{a^{x+h}-a^x}{h}=a^x\cdot\frac{a^h-1}{h}
\]
となる.右辺の $\frac{a^h-1}{h}$ は $x$ を含まないので,この極限,すなわち $x=0$ における $a^x$ の傾きさえ分かればよい.底 $a$ を変えて,\,$h=0.001$ のときの値を計算してみると,
\begin{center}
\begin{tabular}{ccccc}
\toprule
$a$ & $2$ & $2.5$ & $2.7$ & $3$\\
\midrule
$\frac{a^h-1}{h}$ & $0.6934$ & $0.9167$ & $0.9937$ & $1.0992$\\
\bottomrule
\end{tabular}
\end{center}
となる.底を $2$ から $3$ へ大きくすると,値は $1$ より小さい数から $1$ より大きい数へ変わる.その間に,傾きがちょうど $1$ になる底があると予想できるだろう.そのような底を選べば,上の式の係数が $1$ になって話は早い.また,数列 $\left(1+\frac1n\right)^n$ を $n=1,2,10,100,1000$ で計算すると $2,\,2.25,\,2.594,\,2.705,\,2.717$ となり,やはり $2$ と $3$ の間の値に近づいていく.実は,この2つは同じ数なのである.

\begin{theorem}{自然対数の底}
数列 $\left(1+\frac1n\right)^n$ は極限 $\e$ をもち,\,$2<\e<3$ である.この底に対する指数関数は
\[
 \lim_{h\to0}\frac{\e^h-1}{h}=1
\]
を満たす.\,$\e$ を底とする対数を自然対数といい,本書では $\log x$ と書く.
\end{theorem}

この極限を実数指数の定義から構成する議論は,本書の証明範囲から外す.以後の微分公式は,この基本極限と指数法則から導く.\,$\e^{x+h}=\e^x\e^h$ なので,
\[
 \frac{\e^{x+h}-\e^x}{h}
 =\e^x\cdot\frac{\e^h-1}{h}\longrightarrow\e^x.
\]
よって $(\e^x)'=\e^x$ である.

対数では $x>0$ とし,\,$k=\log(x+h)-\log x$ と置く.対数の連続性から $h\to0$ のとき $k\to0$ であり,\,$h\ne0$ なら $k\ne0$ である.また $x+h=x\e^k$ だから,\,$h=x(\e^k-1)$ となる.したがって
\begin{align*}
 \frac{\log(x+h)-\log x}{h}
 &=\frac{k}{x(\e^k-1)}\\
 &=\frac1x\cdot\cfrac{1}{\cfrac{\e^k-1}{k}}\\
 &\longrightarrow\frac1x
\end{align*}
を得る.以上と合成関数の微分から,
\begin{align*}
 (\e^{u(x)})'&=\e^{u(x)}u'(x),\\
 (\log u(x))'&=\frac{u'(x)}{u(x)}\quad(u(x)>0),\\
 (a^x)'&=a^x\log a\quad(a>0)
\end{align*}
となる.最後は $a^x=\e^{x\log a}$ を使った.実験の表の値は $\log a$ の近似値だったのである.

さらに,\,$u(x)<0$ の区間では $\log|u(x)|=\log\{-u(x)\}$ だから,
\[
 (\log|u(x)|)'=\frac{-u'(x)}{-u(x)}=\frac{u'(x)}{u(x)}
\]
となる.したがって,\,$u(x)\ne0$ の区間では $u$ の符号によらず
\[
 (\log|u(x)|)'=\frac{u'(x)}{u(x)}
\]
が成り立つ.絶対値を付けることで,負の値の区間でも実数の対数として定義できる.

$x>0$ では実数 $\alpha$ に対して $x^\alpha=\e^{\alpha\log x}$ と書けるため,
\begin{align*}
 (x^\alpha)'&=\e^{\alpha\log x}\cdot\frac\alpha x\\
 &=\alpha x^{\alpha-1}
\end{align*}
である.整数べきなら,元の関数が定義される負の $x$ にも既に同じ公式が成り立つ.非整数べきの場合は,まず $x>0$ を基本の定義域として考える.

\paragraphbf{対数微分法}
$x^x$ のように底にも指数にも $x$ を含む関数には,\,$x^\alpha$ の公式も $a^x$ の公式もそのままでは使えない.そこで,両辺の対数をとって積の形にしてやる.一般に $f(x)>0$ として $y=f(x)^{g(x)}$ と置くと,\,$y>0$ だから
\[
 \log y=g(x)\log f(x)
\]
と書ける.両辺を $x$ で微分すると,左辺は合成関数の微分から $\frac{y'}{y}$ となるので,
\[
 \frac{y'}{y}=g'(x)\log f(x)+g(x)\frac{f'(x)}{f(x)}
\]
を得る.両辺に $y$ を掛ければ $y'$ が求まる.このように対数をとって積やべきを分けてから微分する方法を\textbf{対数微分法}という.

\begin{bookpoint}
\textbf{POINT}\par
$y>0$ を確かめて両辺の対数をとり,\,$\log y$ を微分して得た $\frac{y'}{y}$ に $y$ を掛ける.
\end{bookpoint}

\begin{example}{微分の計算}
次の関数を微分し,(1) については $x=2$ における接線も求めよ.
\[
 \text{(1)}\quad f(x)=\frac{x^2+1}{x-1},\qquad
 \text{(2)}\quad g(x)=\e^{x^2}\sin(3x),\qquad
 \text{(3)}\quad y=x^x\quad(x>0).
\]
\end{example}
\paragraphbf{方針}
\begin{enumerate}[label=(\arabic*)]
 \item 商の公式を使う.接線は $y-f(2)=f'(2)(x-2)$ に代入する.
 \item 外側の積の公式を使い,各因子は合成関数の微分で微分する.
 \item 対数微分法を使う.
\end{enumerate}
\paragraphbf{解答}
\begin{enumerate}[label=(\arabic*)]
 \item $f'(x)=\frac{x^2-2x-1}{(x-1)^2}$,接線は $y=-x+7$
 \item $g'(x)=2x\e^{x^2}\sin(3x)+3\e^{x^2}\cos(3x)$
 \item $y'=x^x(\log x+1)$
\end{enumerate}
\paragraphbf{解法}
\begin{enumerate}[label=(\arabic*)]
 \item $x\ne1$ で商の公式を使うと,
 \begin{align*}
  f'(x)&=\frac{2x(x-1)-(x^2+1)\cdot1}{(x-1)^2}\\
  &=\frac{x^2-2x-1}{(x-1)^2}
 \end{align*}
 となる.\,$f(2)=5$,\,$f'(2)=-1$ なので,接線は $y-5=-(x-2)$,すなわち $y=-x+7$ である.

 別の計算として,多項式の割り算から
 \[
  f(x)=x+1+\frac2{x-1}
 \]
 と書き直してから微分すると,
 \[
  f'(x)=1-\frac2{(x-1)^2}
 \]
 となり,\,$f'(2)=-1$ が早く求まる.微分する前に式を整理すると,計算が短くなる場合がある.
 \item 2つの関数 $\e^{x^2}$ と $\sin(3x)$ の積であり,それぞれの中に合成がある.まず外側の積の公式を使い,次に内側を微分すると,
 \[
  g'(x)=2x\e^{x^2}\sin(3x)+3\e^{x^2}\cos(3x)
 \]
 となる.\,$2x$ と $3$ は,それぞれ $x^2$ と $3x$ を微分した因子である.
 \item $x>0$ より $y=x^x>0$ だから,両辺の対数をとって
 \[
  \log y=x\log x
 \]
 とできる.両辺を $x$ で微分すると,
 \[
  \frac{y'}y=\log x+1
 \]
 となるので,
 \[
  y'=x^x(\log x+1)
 \]
 である.
\end{enumerate}
